\section*{Abstract}
Multimodal models trained by empirical risk minimization (ERM) can exploit spurious visual correlations in their training data instead of the physiological signal they are intended to use, and this risk is greatest when data are scarce. This thesis introduces EQUITAS-RCMF, a multimodal architecture for stress detection that constrains the structure of the network's visual feature space rather than relying on regularization alone. Its central component is a Stiefel orthogonal subspace decomposition, which projects MobileNetV3-Small visual features onto two mutually orthogonal subspaces, a causal subspace and a confounder subspace, using a thin QR projection at every forward pass. Following the learning-using-privileged-information (LUPI) paradigm, a ground-truth scar mask supervises the confounder subspace during training and is not required at inference, where the scar-related focus signal is estimated from the confounder subspace instead. A gate driven by this focus signal reduces the contribution of visual features when they concentrate on the scar region, and the training objective combines counterfactual consistency with demographic-parity and equalized-odds penalties.

Experiments use a pre-registered pilot cohort of eight participants from the UBFC-Phys social-stress dataset with a synthetic brow-line scar. A pre-registered test shows that the facial video itself contains recoverable physiological information (POS estimator, Mann--Whitney $p = 0.0314$; exploratory CHROM $p = 0.0001$ and PBV $p = 0.0103$), so a visual branch has access to both genuine and spurious cues. In the reported demographic audit, EQUITAS-RCMF shows a counterfactual gap of approximately $0.0006$ between predictions for scarred and scar-removed inputs, and a JIT-compiled inference graph runs at 0.90~ms per frame (1112.7 frames per second) on an x86-64 CPU, well within a 30-frame-per-second real-time budget. The thesis also discusses the limits of these results, including the small cohort, the pending recovery of the master checkpoint, and modality collapse: the fused model is currently less accurate than a physiology-only baseline. It also examines the design's implications for data protection and algorithmic accountability.

\vspace{0.75cm}

\noindent\textbf{Keywords:} Counterfactual Fairness; Edge AI; Stiefel Manifold; Learning Using Privileged Information; MobileNetV3; Remote Photoplethysmography; Algorithmic Bias; Stress Detection
